\chapter{Alien ABDUCTION \& PARAPSYCHOLOGY}\label{ch:abduction}
\markboth{Alien Abduction \& Parapsychology}{Alien Abduction \& Parapsychology}

\chapquote{``In order to understand more it is imperative that we improve our knowledge before choosing which side of the fence we feel compelled to belong.''}{J.\,P.~Robinson}

% ---------------------------------------------------------------
\section{What is Hypnosis?}\label{sec:hypnosis}

Hypnosis is a state of focused attention with heightened suggestibility. That is the whole of
it. It is not sleep --- EEG readings during hypnosis resemble waking attention, not any sleep
stage. It is not unconsciousness, it does not surrender the will, and it cannot make a person
do something they find genuinely abhorrent. \hi{What it reliably does is lower the threshold at
which a suggestion becomes an experience.}

It has real clinical uses. There is decent evidence for hypnosis in pain management, in
irritable bowel syndrome, in smoking cessation as an adjunct, and in anxiety around dental and
surgical procedures. Some people are highly susceptible and some are barely susceptible at
all; it appears to be a stable trait, roughly normally distributed.

What hypnosis absolutely does not do --- and this is not a controversial position, it is the
settled position of the relevant professional bodies --- is recover accurate memory.

Understand why, because the mechanism matters. Memory is not storage. \hi{It is not a recording
that can be replayed with sufficient effort.} Every act of recall is an act of reconstruction,
assembled from fragments, and the reconstruction is rewritten and re-stored each time,
incorporating whatever context is present at the moment of recall. Elizabeth Loftus spent a
career demonstrating this experimentally. Change one word in a question --- ``how fast were
the cars going when they \emph{smashed}'' versus ``\emph{hit}'' --- and you change both the
speed estimate and whether subjects later remember broken glass that was never there. Her
group implanted a false childhood memory of being lost in a shopping mall in roughly a
quarter of subjects, using nothing but a suggestion from a relative.

Now add hypnosis. Under hypnosis the subject is more suggestible, less critical of
implausibility, and --- this is the destructive part --- reports the resulting material with
\emph{increased subjective confidence}. Hypnotic recall produces more detail, more vivid
detail, and more errors, while making the subject certain. Confabulation under hypnosis is
not lying. The person genuinely remembers it afterwards, and cannot distinguish it from a real
memory, because neurologically it now is one.

\begin{funfact}
The legal system worked this out and acted on it. Hypnotically refreshed testimony is
inadmissible or severely restricted in most US jurisdictions, and the American Psychological
Association, the American Medical Association and the British Psychological Society have all
issued statements against its use for memory retrieval. The reason was not theoretical: it was
the wave of wrongful convictions and the ``satanic ritual abuse'' panic of the 1980s and
1990s, in which therapists using regression techniques generated thousands of detailed,
confidently-held, entirely false memories, destroyed families, and imprisoned innocent people.
The abduction literature was built with the same tools, in the same decade, frequently by
people with no clinical training at all.
\end{funfact}

So here is my position, stated plainly, because it governs the rest of this chapter. \hi{Any
account whose content originates under hypnosis is not evidence about the world.} It is
evidence about the session. This is not a claim that the witnesses are dishonest --- almost
none of them are, and many are clearly suffering. It is a claim about what the instrument can
measure. Handing a sincere frightened person to a regression practitioner with a theory does
not recover what happened to them. It replaces it.

% ---------------------------------------------------------------
\section{Betty and Barney Hill}\label{sec:bettybarneyhill}

On the night of 19 September 1961, driving home to Portsmouth, New Hampshire on Route 3, Betty
and Barney Hill saw a light in the sky that appeared to follow their car. Barney stopped,
looked at it through binoculars, and became frightened. They drove home in a hurry, arriving
about two hours later than they expected. This is the founding case of the entire abduction
genre --- everything that follows in the literature, from the medical examination to the
missing time to the greys themselves, begins here.

What makes the case invaluable is that its development is fully documented, step by step, which
is almost unique. We can watch the story being built.

In the days immediately afterwards, the Hills reported a sighting. That is all. They reported
it to Pease Air Force Base, and the Blue Book entry describes lights and a possible aircraft.
Betty then borrowed a book by Donald Keyhoe from the library. Within about ten days she began
having a series of vivid nightmares involving being taken aboard a craft, examined by
humanoid beings with large eyes, and having a needle inserted into her navel. She wrote the
dreams down. She discussed them with Barney, repeatedly, over the following months.

\figwrap[l]{0.38}{Betty and Barney Hill, whose 1961 encounter became the template for the abduction narrative.}{https://commons.wikimedia.org/wiki/File:Betty_and_Barney_Hill.jpg}{fig:hills}
Two and a half years later, in early 1964, both Hills underwent hypnotic regression with Dr.~Benjamin Simon (see Figure~\ref{fig:hills}), a Boston psychiatrist. Under hypnosis they produced detailed abduction
accounts. Barney's account included elements matching Betty's dreams, which he had by then
heard many times. Simon --- and this is the part almost never quoted --- concluded that the
abduction narrative was not a real memory. His professional assessment was that it was a
fantasy, originating with Betty's dreams, transferred to Barney through their conversations,
and that the Hills sincerely believed it. He said so publicly and repeatedly for the rest of
his life.

The case reached the public through a 1965 newspaper leak and then John Fuller's 1966
bestseller \emph{The Interrupted Journey}, followed by a 1975 television film watched by tens
of millions. The greys entered popular culture through that broadcast.

The Hill case is not remarkable because it is false. The Hill case is remarkable. Not
because of the claims of alien abduction. It is remarkable because a woman managed to
produce an accurate star map years before the scientific discovery of those stars solidified
the structure of that star map.

Betty, under hypnosis, said she was shown a map of trade routes and later drew it from
memory. In 1968 an amateur astronomer, Marjorie Fish, spent years building physical models
of nearby star systems and found an orientation matching Betty's sketch, with Zeta Reticuli
as the origin. This was published in \emph{Astronomy} magazine in 1974 and caused a
sensation.

Fish's specific Zeta Reticuli identification does not survive the modern data. The match
depended on the stellar distance data available in the early 1970s; when the Hipparcos
satellite produced accurate parallax measurements in the 1990s, two of the central stars in
Fish's model moved catastrophically --- Kappa Fornacis shifted from 42 to 105 light-years,
HD~13435 from 37 to 188 --- and the ring she identified no longer holds together. Carl
Sagan and Steven Soter also pointed out at the time that the exercise was statistically
empty: with enough stars to choose from and freedom to select the viewing angle, the scale,
and which dots count, you can match almost any arrangement of dots. Take Betty's sketch,
they noted, and you can find matches to random dot patterns. This is the
Section~\ref{sec:leylines} ley-line problem in three dimensions, and it is a real
objection.

But the statistical objection and the Hipparcos correction are not the same thing, and
they pull in opposite directions. The statistical objection says the match is empty
because anything fits. The Hipparcos correction says the match is wrong because the wrong
stars were used. Both cannot be true in the same sense at the same time. If the pattern is
so flexible that anything matches, then the Hipparcos corrections should have produced
another equally convincing match with the updated catalogue --- and in 2026 a recomputation
using full Gaia and Hipparcos data found exactly that: the pattern can indeed be reproduced
from other nearby stars, just not from the ones Fish named. The hub of the arrangement
shifted to two cool dwarfs in Camelopardalis, about as far from Reticulum as you can get on
the sky. The match was tight. Random dot patterns fit just as well. Twenty-one anonymous dots
without brightnesses, colours or scale simply under-determine a three-dimensional viewpoint
among thousands of candidates. The information is not there.

And yet. Betty Hill was not an astronomer. She had no star catalogue, no fishing-line model,
no five-year search. She drew twenty-one dots from a hypnotic memory three years after the
fact, and those dots matched real stars closely enough that a dedicated amateur spent half a
decade confirming it and a major astronomy magazine published the result. The information
may not have been enough to locate a viewpoint, but it was enough to look like one. That is
what makes the case remarkable.

\actualq{Did it happen?}{How was the star map produced so accurately?}

\begin{funfact}
There is a strong candidate for what the Hills actually saw, and it is oddly poignant. On that
night Jupiter and Saturn were both prominent low in the southwest along their route --- a
bright object and a fainter companion, which is exactly the initial description --- and a
vehicle driving a winding road through the White Mountains will see such objects apparently
move, stop, and follow. Barney was an interracial-marriage-era postal worker in 1961 New
Hampshire, active in the NAACP, driving at night through rural country, already anxious. The
emotional weight of that night was real. What it was about may have had less to do with the
sky than with the road.
\end{funfact}

% ---------------------------------------------------------------
\section{Fire in the Sky}\label{sec:fireinthesky}

On the evening of 5 November 1975 a seven-man logging crew was thinning timber at Turkey
Springs, in the Apache-Sitgreaves National Forest about twelve miles south of Heber, Arizona.
Driving out at dusk they saw a glowing object off the road. Travis Walton, aged twenty-two,
got out of the truck and walked towards it. The other six men said a beam of light struck him,
and that they panicked and drove away; when they came back he was gone. They reported him
missing to the Navajo County sheriff's office that same night, around a quarter to eight. A
search involving horses, a helicopter and a Geiger counter turned up nothing at all --- no
tracks, no burns, no clothing, no body. Just after midnight on 11 November, five days and six
hours after he was last seen, Walton rang his sister's house collect from a payphone in Heber.

This is the case the abduction literature reaches for when it wants to be taken seriously, and
I understand why. Six independent witnesses reported the event before there was anything to
report; five of them passed a polygraph five days later; the family had no obvious way of
knowing the story would end well. It has a film, \emph{Fire in the Sky}, discussed in
Chapter~\ref{ch:media}, and a book, and half a century of retelling. It is also, in my view,
the most instructive case in this chapter, because everyone argues about it in the wrong
currency.

\begin{funfact}[Fun Fact: The Author Met Him]
In the summer of 2024 I got the chance to meet Travis Walton himself in Roswell, New Mexico, and
he signed my \emph{Fire in the Sky} poster. I want to say plainly where I stand, because readers
are entitled to know it: I believe Travis Walton. I do not disagree with his account of his own
life. He was generous with his time, and he has told the same story, under pressure, for half a
century, at considerable personal cost and with nothing much to show for it. What follows in the
next few pages is not a case against the man. It is the one physical question inside the case
that nobody on either side has answered, and I would rather put it to him honestly than pretend
it is not there.\par
\vspace{7pt}
\begin{center}
\adjustimage{max size={0.60\linewidth}{3.5in}}{walton_autograph.jpg}\par\vspace{4pt}
\begin{minipage}{0.86\linewidth}
{\RaggedRight\scriptsize\color{steel} \emph{Fire in the Sky} (Paramount, 1993) one-sheet, signed
by Travis Walton in Roswell, New Mexico, summer 2024: ``Evan, Keep an eye on the sky --- and your
feet on the ground!'' Author's own photograph of his own copy.\par}
\end{minipage}
\end{center}
\end{funfact}

The standard debate runs: did aliens take him, or did the crew stage a hoax? Both sides then
fight over motive, character and lie-detector charts, none of which settle anything, because a
polygraph measures anxiety and a motive is not a mechanism. Meanwhile the case contains one
question with a hard physiological answer, and almost nobody asks it.

\actualq{whether Travis Walton was abducted.}{what he drank. A man goes without water for
about three days before it starts killing him --- so who gave Travis Walton a cup of water for
five days and six hours?}

Here is the problem. \hi{Wherever Walton was, he was somewhere with food and water.} The human
limit without water is roughly three days --- the old wilderness rule of threes --- and while
cool air and rest can stretch that towards five or six, nobody arrives at the far end of it
looking merely dishevelled. Severe dehydration is not subtle. You get tachycardia, low blood
pressure, sunken eyes, cracked mucous membranes, dark and drastically concentrated urine,
raised serum sodium, and confusion. A man who had also eaten nothing for five days would be
running on fat, and a body running on fat dumps ketones into the urine, which is trivially easy
to detect and was standard practice in 1975.

Walton was examined on 11 November by Dr.~Joseph Saults and Dr.~Howard Kandell. He was thirsty
and said he had lost about ten pounds. What he was not was clinically dehydrated, and the
urinalysis found no meaningful ketones. \hi{That single result is the whole case.} A man five
days into total starvation has ketonuria; a man who ate a few hours ago does not. So he ate. And
if he ate, the five days were not spent unconscious on an examination table --- they were spent
somewhere with a kitchen, a cool box, or at minimum a canteen and a tin.

\subsection{Did They Carry Water?}

There is a further point, and the abduction reading rarely confronts it. If Walton was fed and
watered aboard the craft, then his hosts were carrying provisions metabolically compatible with
a human being --- water at our tolerances of salinity and temperature, calories in a form a
primate gut can absorb --- and that is a far larger claim than a light in the sky. A crew that
has crossed interstellar distance has no obvious reason to stock a larder for a species it has
not yet met, and if it did, we have left the realm of anomaly altogether and entered that of a
logistics operation with a menu.

I want to take that seriously rather than use it as a punchline, because it is the one part of
the abduction reading that can be costed. Start with the water. A resting adult male loses
something like two to two and a half litres a day through urine, breath and skin, and he cannot
bank it in advance. Five days and six hours is therefore roughly twelve litres, and twelve
litres is twelve kilogrammes --- more mass than the man's own head. Add the food. Even at a
sedentary floor of around fifteen hundred kilocalories a day he needs some three to four
kilogrammes of anything worth eating. \hi{So a crew holding one human being for five days is
carrying about sixteen kilogrammes of consumables for him alone.} That is not an exotic
objection. It is the same arithmetic that governs every crewed vehicle we have ever flown, and
mass is the one currency no propulsion system makes free.

Then there is the chemistry, which is worse. Water is water anywhere in the galaxy, but potable
water is a narrow specification: not brackish, not heavy in deuterium, not carrying dissolved
metals, not so cold it triggers gastric cramp in a shocked man, and not sterile in the wrong
direction either. Food is narrower still. Every protein in your body is built from left-handed
amino acids and every sugar you can metabolise is right-handed, and that handedness is a local
accident of terrestrial biochemistry, not a law of the universe. A biosphere that came out the
other way round would produce food that is nutritionally inert to us --- edible in the sense
that sand is edible. Which leaves a crew with exactly two options: synthesise glucose, salts
and water to human specification, or carry human food. \hi{Both options require them to have
studied human metabolism before they picked him up.} And if they had already studied it, Walton
was not a discovery. He was a repeat customer.

Now put the man himself back in the picture, because the abduction literature quietly assumes
he was unconscious for most of those five days, and an unconscious body is much harder to keep
than a conscious one. You cannot pour water into a sedated man; he aspirates it. You need a
nasogastric tube or an intravenous line, and both leave marks. Five days of continuous fluids
means repeated cannulation --- typically a fresh site every two or three days --- so several
puncture wounds at various stages of healing, bruising around them, and very probably a urinary
catheter as well. Five days flat on a table also means the beginnings of pressure damage over
the sacrum and heels, measurable muscle wasting, and the sort of gut that has not moved in a
week. Drs.~Saults and Kandell found none of it. What they found was one red mark on the inside
of the right elbow, which Kandell described as compatible with a puncture such as when somebody
takes blood from you, and which was judged to be a day or two old --- not five. Duane had
already noticed that mark before either doctor arrived, and so had the hypnotherapist. A single
fresh needle mark is what you get from a blood draw, or from an injection, or from a scratch on
a fence wire. It is not the record of a week on a drip.

So the provisioning question folds back into the metabolic one and comes out the same way. If he
was conscious and fed, somebody handed him a cup and a plate, and that somebody had to know
what a human being drinks. If he was unconscious and fed, his body would carry a fortnight's
worth of clinical evidence, and it carried a graze on one elbow. If he was neither fed nor
watered, his urine would have been dark and loaded with acetone, and it was clear and had none.
The ten pounds he said he had lost fits a few uncomfortable days of light eating, stress and
broken sleep. It does not fit total starvation, which costs a man closer to a stone and leaves
signs on his tongue, his skin and his blood pressure that two doctors examining him in a
living room could not have missed.

I should be fair about the escape hatch, since one exists. A civilisation capable of crossing
interstellar space is presumably capable of distilling water and running sugars out of a
synthesiser, and could in principle feed a captive perfectly, painlessly and invisibly. I
cannot rule that out. What I can point out is the price of the manoeuvre: it converts a
surprise encounter on a logging road into a planned operation, staffed by people who had
already worked out our salt tolerance and our chirality, and who then dropped their subject at
a payphone in Heber. \hi{Every version of the abduction that survives the urine sample requires
the visitors to have been catering for us in advance.} A cool box in the back of a pickup, or a
fire lookout's kitchen five miles up the road, requires nothing at all.

Notice that this does not require me to call anyone a liar about the object by the road. It
simply removes the missing five days from the realm of the unexplained and puts them in the
realm of the ordinary. And it is worth pointing out that Walton's own account does not cover
those days either. He describes something in the order of two hours --- a table, thin figures,
a helmeted human-looking man, a hangar --- and then a road outside Heber. Even on his own
telling, roughly five days are simply absent. \hi{The abduction narrative is not an explanation
of the missing time; it is a story attached to the beginning and the end of it.}

\subsection{The Terrestrial Reading}

Once the metabolic question is settled, the terrestrial reading of the case stops being a
sneer and starts being merely economical. The crew's foreman, Mike Rogers, held a US Forest
Service thinning contract on 1,277 acres that was badly behind; an inspector had told him in
mid-October that the remaining work could not be done in the time left; he had obtained an
extension to 10 November, and the contract carried a penalty of roughly \$2,500 --- serious
money in 1975 --- with the usual clause excusing an act of God. The night Rogers wrote to the
Forest Service about his difficulties, 20 October 1975, NBC broadcast \emph{The UFO Incident},
the dramatisation of the Hill case discussed in Section~\ref{sec:bettybarneyhill}. Two weeks
later the crew reported that one of their number had been taken by a flying saucer, and the
work stopped. The \emph{National Enquirer} was at that time offering a \$5,000 prize for the
best UFO case of the year, and paid it to them.

The polygraph evidence is weaker than it sounds for the same structural reason. The crew were
asked, in substance, whether they had harmed Walton, whether they knew where his body was, and
whether they had seen a UFO. If two men staged something for the benefit of the other five,
every man in the truck can answer all four questions honestly. Walton himself was given a
polygraph by the \emph{Enquirer}'s examiner on his return and failed it badly; that result was
kept quiet for twenty years, and a subsequent test scheduled at Holbrook was not attended. The
later tests he passed were arranged by sympathetic parties, and by then he had had years of
practice at telling the story. He had also, in 1971, pleaded guilty to forging cheques --- which
proves nothing about 1975, but does bear on the recurring claim that no one involved had any
history of dishonesty.

As for where he actually was, the most economical suggestion I have seen comes from the
Australian researcher Charlie Wiser and the producer Ryan Gordon, working partly from
conversations with the crew: the Gentry fire lookout tower, five miles by road from the work
site, seventy feet up, with a lit fourteen-foot-square live-in cab that was never searched
because the sheriff was running a missing-person inquiry and had no reason to ask the lookout
what he had seen. Steve Pierce, the youngest of the crew and now the most candid, has said his
uncle drove him out to the tower while Walton was still missing and told him plainly that this
was the UFO. Rogers, recorded in 2021, appeared to describe the whole thing as staged, then
retracted the remark, then partially retracted the retraction. I do not present any of that as
proof. I present it as the sort of mundane possibility that a search party never checked, and
that a fed, hydrated, ketone-free man returning to a payphone in Heber makes unnecessary to
rule out.

\begin{funfact}
Two details about that tower are hard to unsee once noticed. The first is the shape: the
earliest witness sketches, made within a fortnight of the event, show a squat, angular,
window-framed object roughly fifteen to twenty feet across and eight to ten feet thick --- a
reasonable description of a lookout cab lit from inside, and a poor description of the sleek
silver disc that appeared in later paintings and in the 1993 film. The second is the shift
pattern. Forest Service lookouts staffed those cabs in rotations of five days.
\end{funfact}

I want to be careful about what I am claiming. I do not know what happened at Turkey Springs on
5 November 1975, and I am not in a position to prove a hoax. What I will say is that the case is
not the impasse it is usually presented as. \hi{It has one piece of hard physical evidence in it,
and that evidence is a urine sample.} Everything else --- the sketches, the charts, the
hypnotherapist Walton saw on 11 November before any physician examined him, which by
Section~\ref{sec:hypnosis} is roughly the worst possible order of operations --- is testimony
about testimony. The laboratory result is not. A body keeps its own records, it cannot be
coached, and it says he was eating.

% ---------------------------------------------------------------
\section{Close Encounters}\label{sec:closeencounters}

Hynek's classification, introduced in \emph{The UFO Experience} in 1972, is the field's one
genuinely useful piece of taxonomy, and I gave you the outline in Section~\ref{sec:classificationcommonufos}. Here is what
each level is actually worth.

\textbf{Close Encounter of the First Kind.} An object observed at close range --- Hynek used
about 500 feet --- with no interaction and no physical trace. Evidentially this is a
sighting with better resolution. It rises or falls entirely on witness quality and on what
Section~\ref{sec:investigativetechniques} step two turns up about the sky.

\textbf{Close Encounter of the Second Kind.} An encounter leaving physical effects: scorched
or depressed ground, damaged vegetation, vehicle engine or electrical failure, radiation
readings, animal distress, burns. This is the only category with real evidentiary potential,
because physical traces can be sampled, controlled, and analysed by people who do not know
where the sample came from. The Ted Phillips catalogue collected several thousand physical
trace cases. The frustrating truth is that the great majority of them were sampled without
controls, which as Section~\ref{sec:investigativetechniques} explained renders the laboratory results meaningless. The
best-known CE2 in the field, Cash--Landrum, appears in Section~\ref{sec:cashlandrum}, and it
is the exception that shows what the category could be.

The other case worth putting beside it is Falcon Lake, in Manitoba, on 20 May 1967, because it
is the closest thing this field has to a CE2 with a paper trail: a lone prospector, Stefan
Michalak, burned in a geometric grid pattern across his chest by what he described as a blast of
hot gas vented from the side of a landed craft, with a scorched shirt, a melted glove, a hospital
record, and files opened by the Royal Canadian Mounted Police, the Royal Canadian Air Force and
the Canadian Department of National Defence. It is the one case in this book where the injury
itself is the evidence, and because the burn mechanism he described is an exhaust-gas burn rather
than a radiation burn, I have given it a full treatment alongside the other physical-trace and
radiation-claim cases in Section~\ref{sec:falconlake}. Read the two together. They are the shape
of the best evidence ufology has, and you should get used to it.

\textbf{Close Encounter of the Third Kind.} Occupants observed. Hynek was reluctant to include
this category and said so, on the grounds that it would damage the credibility of the rest of
the scale. He included it because the reports existed and he refused to discard data he found
embarrassing --- which is, for my money, the single most admirable act in the history of this
field.

\textbf{Close Encounter of the Fourth Kind.} Abduction. Added later, not by Hynek. Note that
the scale's evidential quality is now inversely proportional to its number: CE2 has physical
samples, CE4 has a hypnosis transcript.

\textbf{Close Encounter of the Fifth Kind.} Deliberate human-initiated contact, promoted from
the 1990s by advocates of ``CE5 protocols'' --- group meditation to summon craft. There is no
controlled evidence for any of it, the reported results are precisely what expectation and
group dynamics would produce in a field at night, and the practice has no place on a scale
intended to classify observations.

\subsubsection*{The two levels the internet added}

Beyond the fifth kind the scale stops being taxonomy and becomes folklore, and it is worth being
explicit about where it came from. \hi{Neither of the next two categories is Hynek's, and neither
has a traceable primary source.} They were assembled on the open web --- forum posts, wiki edits,
list articles and later video essays --- during the 2000s and 2010s, propagated by repetition
rather than by publication, and they appear in no peer-reviewed paper and in none of Hynek's
work. Hynek died in 1986. He proposed three close-encounter levels in \emph{The UFO Experience}
in 1972 and nothing beyond them. I include the two additions here because you will meet them
constantly and should know what they are, not because they carry any standing.

\textbf{Close Encounter of the Sixth Kind.} An encounter in which an extraterrestrial entity
causes direct physical injury or death to a human being. As a definition this is at least
checkable in principle, which makes it the more interesting of the two: injury leaves a
record, and death leaves a coroner. That is also precisely the problem with it. Every case
advanced as a CE6 either has the injury without the entity --- Cash--Landrum in
Section~\ref{sec:cashlandrum} and Falcon Lake in Section~\ref{sec:falconlake} are both
ordinary CE2s with a medical file attached, and neither witness saw who or what was aboard ---
or has the entity without a body, as with the animal-mutilation and alleged
military-casualty stories where no remains, no autopsy and no death certificate are ever
produced. The category asks for a homicide with an off-world suspect, and the field has not
yet produced a single instance with both a victim and an attacker.

\textbf{Close Encounter of the Seventh Kind.} A purported involvement of human--alien
hybridisation, up to and including the production of a hybrid child. This is the terminal point
of the abduction narrative described below, and it is the least testable proposition on the
scale despite being the one most easily settled: a hybrid child is a genome, and a genome can be
sequenced for a few hundred dollars in any commercial laboratory. No claimant has ever produced
tissue. \hi{The category makes the one claim in ufology that a single cheek swab could confirm
tomorrow, and in forty years not one swab has been offered.} Chapter~\ref{ch:life} deals with
why cross-species reproduction between independently evolved biospheres is not merely unlikely
but structurally impossible --- separate chromosome counts, incompatible genetic coding, no
shared ancestry to hybridise from --- so I will not repeat the argument here. Note only what the
placement of these two levels does to the scale as a whole: it now ends with its two least
evidenced claims sitting at the top, which is exactly backwards, and is the clearest
illustration I know that the number of a close encounter tells you about the size of the story
and nothing whatever about the strength of the evidence.

\Needspace*{3.4in}
% ---------------------------------------------------------------
% fig_cescale.tex --- the close encounter scale as a left-to-right
% progression, for Section~\ref{sec:closeencounters}. Drawn in TikZ so it
% prints sharp at any size and carries no image credit.
% Colour coding: green = Hynek's own three levels, amber = the two later
% additions, red = the two internet-era additions.
% ---------------------------------------------------------------
\begingroup
\sffamily
\begin{center}
\refstepcounter{figph}\label{fig:cescale}%
% Drawn at a generous scale and squeezed to the measure on the way in, so the
% chevron type lands at roughly 7pt on the printed page.
\resizebox{\linewidth}{!}{%
\begin{tikzpicture}[x=1cm,y=1cm,
  cesig/.style={signal,signal to=east,signal pointer angle=140,
                draw=ufogreenink,line width=0.5pt,minimum width=3.0cm,
                minimum height=1.6cm,inner sep=1pt,align=center,
                text width=2.25cm,text=ufogreenink,
                font=\fontsize{8.4}{9.6}\selectfont},
  band/.style={draw=none,inner sep=0pt,anchor=south west},
  bandtext/.style={anchor=north,text=steel,align=center,
                font=\fontsize{8.2}{9.4}\selectfont},
  axistext/.style={text=steel,font=\fontsize{8.6}{9.8}\selectfont}]

% ---- direction of travel, above the row ----
\draw[->,line width=0.9pt,ufogreendark] (-1.5,1.62) -- (17.76,1.62);
\node[axistext,anchor=south] at (8.13,1.78)
  {increasing claimed intimacy of contact \;---\; and, as a rule, decreasing evidential quality};

% ---- the seven levels ----
\node[cesig,signal from=nowhere,fill=ufogreen!55] at (0.00,0)
  {\textbf{CE1}\\[-0.2ex]{\fontsize{7.6}{8.6}\selectfont close sighting}};
\node[cesig,signal from=west,fill=ufogreen!40] at (2.71,0)
  {\textbf{CE2}\\[-0.2ex]{\fontsize{7.6}{8.6}\selectfont physical trace}};
\node[cesig,signal from=west,fill=ufogreen!26] at (5.42,0)
  {\textbf{CE3}\\[-0.2ex]{\fontsize{7.6}{8.6}\selectfont occupants seen}};
\node[cesig,signal from=west,fill=amber!30] at (8.13,0)
  {\textbf{CE4}\\[-0.2ex]{\fontsize{7.6}{8.6}\selectfont abduction}};
\node[cesig,signal from=west,fill=amber!18] at (10.84,0)
  {\textbf{CE5}\\[-0.2ex]{\fontsize{7.6}{8.6}\selectfont contact invited}};
\node[cesig,signal from=west,fill=gapred!26] at (13.55,0)
  {\textbf{CE6}\\[-0.2ex]{\fontsize{7.6}{8.6}\selectfont injury or death}};
\node[cesig,signal from=west,fill=gapred!42] at (16.26,0)
  {\textbf{CE7}\\[-0.2ex]{\fontsize{7.6}{8.6}\selectfont hybrid child}};

% ---- provenance bands beneath ----
\fill[ufogreen!55]  (-1.50,-1.42) rectangle (6.775,-1.10);
\fill[amber!30]     (6.775,-1.42) rectangle (12.195,-1.10);
\fill[gapred!34]    (12.195,-1.42) rectangle (17.76,-1.10);
\draw[midgrey,line width=0.3pt] (-1.50,-1.42) rectangle (17.76,-1.10);
\draw[midgrey,line width=0.3pt] (6.775,-1.42) -- (6.775,-1.10);
\draw[midgrey,line width=0.3pt] (12.195,-1.42) -- (12.195,-1.10);
\node[bandtext] at (2.64,-1.54) {Hynek's original scale,\\\emph{The UFO Experience}, 1972};
\node[bandtext] at (9.49,-1.54) {added later,\\not by Hynek};
\node[bandtext] at (14.98,-1.54) {internet-era additions,\\no formal standing};
\end{tikzpicture}}\par\vspace{4pt}
{\footnotesize\color{steel}\textbf{\textcolor{ufogreenink}{Figure \thefigph.}} The close
encounter scale read left to right. Only the first three levels are Hynek's; the fourth and
fifth were bolted on by other hands, and the sixth and seventh were assembled on the open web
without a traceable primary source. The evidential value of the scale runs in the opposite
direction to its numbering.}
\end{center}
\endgroup


The abduction narrative itself has a remarkably fixed structure, and structure is diagnostic.
The sequence is: capture (often from a bedroom, often preceded by paralysis), transit through
a beam or corridor, examination on a table by beings performing medical procedures, focus on
reproductive organs, a ``conference'' or telepathic communication frequently including an
environmental warning, a tour of the craft, return, and missing time. Budd Hopkins and David
Jacobs treated this consistency as corroboration --- how could thousands of unconnected people
report the same thing?

The answer is that they were not unconnected. \emph{The Interrupted Journey} was a bestseller
in 1966; the television film reached tens of millions in 1975; \emph{Communion} sold in the
millions in 1987 with a grey face on the cover. By the time the great mass of cases arrived,
the script was in every supermarket in America. And the majority of those cases were developed
by a small number of investigators using regression, asking questions shaped by the accounts
they had already collected. Section~\ref{sec:confirmationbias} and Section~\ref{sec:hypnosis} between them predict
exactly this outcome: high consistency, high confidence, no independent corroboration.

\subsubsection*{Budd Hopkins, and the case that broke his method}

Because Hopkins's name has now appeared several times, he deserves a proper account rather than
a label. Elliot Budd Hopkins (1931--2011) was not a scientist and never claimed to be one. He
was an abstract expressionist painter of genuine standing, with work in the Whitney, the
Metropolitan Museum of Art and the Hirshhorn, who saw a daylight disc in 1964 and spent the
remaining forty-seven years of his life on the subject. \emph{Missing Time} (1981) introduced
the idea that an abduction could be concealed from the experiencer entirely, recoverable only
under hypnosis; \emph{Intruders} (1987) added the breeding programme; and in 1989 he founded
the Intruders Foundation in Manhattan to support experiencers. He regressed something on the
order of a thousand people over his career.

\figwrap[l]{0.34}{Budd Hopkins addressing UFO investigators at a New York conference. A serious,
decent man, entirely self-taught in the one discipline his conclusions depended on.}{}{fig:hopkins}
I want to be careful here, because it is easy and cheap to be contemptuous of him. Everyone who
dealt with Hopkins, including his critics, described the same person: patient, courteous, wholly
sincere, and unusually kind to frightened people at a time when no institution in the United
States would give them ten minutes. He took no money from the experiencers he worked with. He
kept records. Several professional psychiatrists, John Mack among them, came to the subject
because Hopkins persuaded them the distress was real. That much is to his credit and it is
permanent.

\setlength{\figwraprise}{\baselineskip}%
\figwrap[r]{0.34}{Hopkins conducting a hypnosis session with an experiencer. This is the
instrument that produced most of the abduction corpus --- an unlicensed practitioner, a
believed-in hypothesis, and a suggestible subject in a private room.}{}{fig:hopkinshyp}
The problem is the instrument, and it is the same problem Section~\ref{sec:hypnosis} laid out.
Hopkins had no training in hypnosis, no training in clinical interviewing, and no training in
memory. He believed the breeding-programme hypothesis before he collected most of the data that
supposedly established it, and he conducted the regressions himself, alone, on subjects who
knew what he believed. There is no arrangement more reliably productive of false memory than
that one. It is not fraud. It is the wrong tool, applied thousands of times, with total
conviction.

The case that shows it is the Manhattan abduction of November 1989, which Hopkins regarded as
the most important in the field and published in 1996 as \emph{Witnessed}. His witness, Linda
Napolitano --- given the pseudonym ``Linda Cortile'' --- was said to have been floated out of
the closed window of her twelfth-storey apartment on Manhattan's Lower East Side, in the small
hours, into a hovering craft, in full view of the city. Hopkins eventually claimed some
twenty-three witnesses, including, in the most extraordinary version, a United Nations Secretary
General whose motorcade had allegedly stopped on the FDR Drive to watch.

Consider what that claim requires. A woman is levitated through a sealed window above one of
the most photographed, most densely populated, most sleepless square miles on Earth, and the
entire evidentiary residue is a set of letters sent to Budd Hopkins. No photograph. No
newspaper report the following morning. No police call. No independently traced witness --- the
principal corroborating letters came from correspondents whose identities Hopkins could not
establish, and the investigators Carol Rainey (Jacobs's former wife, who filmed inside the work
for years) and Joe Nickell separately documented how much of the corroboration traced back to
Napolitano herself. Rainey, who had a camera in the room, ultimately published her account of
what the method looked like from the inside, and it is not flattering.

A 2024 Netflix documentary revived the case for a new audience without resolving any of that,
which is why it belongs in this book. Manhattan is the strongest test the abduction hypothesis
has ever been handed --- maximum witnesses, maximum surveillance, maximum documentary
opportunity --- and it produced nothing that survives contact with
Section~\ref{sec:investigativetechniques}. Hopkins deserves to be remembered for taking
frightened people seriously. He also has to be remembered for showing, in the one case where
the city itself was the control group, that the method does not work.

% ---------------------------------------------------------------
\section{Genetic Motivation}\label{sec:geneticmotivation}

The standard explanation offered within the abduction literature for the reproductive focus of
the examinations is a breeding programme: the visitors are harvesting gametes, producing
hybrids, and monitoring the results across generations. Jacobs built an entire late career on
it. It deserves examination on the merits, because it is a biological claim and biology can
adjudicate it.

It does not work, and it fails on grounds that have nothing to do with scepticism about UFOs.

Interbreeding requires close common ancestry. Genetic compatibility is not a matter of
similar appearance --- it requires compatible chromosome number and structure, compatible
gene regulation, compatible developmental timing, and compatible molecular machinery down to
the codon assignments. Humans cannot interbreed with chimpanzees, and we share about 98.8\% of
our DNA with them and a common ancestor six million years ago. We cannot interbreed with any
other species on this planet, all of which descend from the same origin and use the same
genetic code. The proposition that an organism from an independent biosphere --- with its own
chemistry, its own information storage, possibly not using DNA at all --- could produce
viable offspring with us is not merely improbable. It is a category error. It is like
expecting to run Windows software on a plant.

Second, the technology assumption is incoherent. A civilisation capable of interstellar travel
is, necessarily, vastly ahead of us in every technical domain, including molecular biology.
Such a civilisation would not need to abduct farmers repeatedly over decades to obtain genetic
material. A single tissue sample, obtained once, contains the entire genome. We can already
sequence a human genome from a cheek swab for a few hundred dollars, and synthesise
arbitrary DNA sequences to order. A species that crosses light years but requires a
sixty-year sampling programme with a table and instruments is doing molecular biology worse
than we currently do it.

Third, the stated purposes are contradictory across the literature: they need our DNA because
they are dying out; they are creating a hybrid species to replace us; they are correcting a
genetic defect in themselves; they are preparing a vessel for consciousness transfer. Each is
offered with equal confidence by different authors reading the same case files.

And there is a fourth objection, which is the one this book has been assembling in the background
for several chapters and which has not yet been stated outright. Argue about chromosome counts
long enough and you concede the framing --- that the interesting thing here is a cross between two
biospheres. It is not. Look at what is supposedly doing the crossing. The reptilians of
Section~\ref{sec:racesaz} are crocodile hide and snake eyes; the Draconians add a bat's wings and
a rhinoceros horn; elsewhere on the same list sit lion faces, mantis heads, fish skin, insect
limbs and owl eyes. \hi{Every one of them is assembled out of animals from this planet.} An
independent biosphere, evolved under another star, has arrived at a menagerie of Earth's own
phyla --- and it has done so in exactly the manner Section~\ref{sec:animalheads} watched Egyptian
artists work and Section~\ref{sec:animalpeople} watched the myth-makers work, by putting a local
animal's head on an upright human frame. That is not a fact about extraterrestrial genetics. It is
a fact about human image-making, and it is visible before a single chromosome is counted.

\actualq{whether extraterrestrial reptiles are interbreeding with human beings.}{why the aliens
are all hybrids of \emph{Earth-based} species --- why a life-form from another biosphere should
turn out to be built from the animals outside our own window.}

The distinction is worth holding onto, because the two questions have different fates. The
interbreeding question can be argued forever, since neither side can produce the genome. The
Earth-stock question is already answered, by the material itself, and the answer does not require
anybody's cooperation: the catalogue is terrestrial because the people describing it are. The
formal treatment of that answer --- the full chain from animals to anthropomorphic characters to
aliens to hybrids, and the one entry in the taxonomy that breaks it --- is set out at the end of
this book, where the reasoning can be laid beside the evidence rather than asserted here.

\begin{funfact}
There is a much older version of this story, and it is nearly universal. Incubus and succubus
legends across medieval Europe describe nocturnal visitation by entities that paralyse the
sleeper and extract or impart seed. The changeling folklore of Northern Europe describes
children taken and replaced, and fairy abduction describes missing time --- days spent in the
fairy mound that were years outside. Jacques Vall\'ee made this argument in
\emph{Passport to Magonia} in 1969 and it has never been answered: the abduction narrative is
the fairy narrative with the technology updated. Either the phenomenon has changed costume, or
something about human beings produces this story reliably in every century.
\end{funfact}

Which brings me to the explanation I find far more likely, and which is also considerably more
respectful of the witnesses than pretending to take the hybrid programme seriously. Sleep
paralysis is a documented neurological event experienced by something like 8 to 30 per cent of
people at least once. During REM sleep the body's voluntary muscles are inhibited. If
consciousness returns before that inhibition lifts, the person is awake, aware, and completely
unable to move. The state is frequently accompanied by hypnagogic hallucination --- a sensed
presence in the room, figures at the bedside, pressure on the chest, buzzing or humming
sounds, a sense of being watched, floating, and overwhelming terror. Every element of the
classic bedroom abduction onset is a documented feature of a common sleep disorder. Add a
regression session with a practitioner who believes in a breeding programme, and the rest
writes itself.

\subsubsection*{The hybrids who came forward: Charmaine D'Rozario-Saytch}

The breeding-programme claim has an end state that is rarely examined, and examining it is
instructive. If hybrids were produced, some of them should now be adults, living among us, and
able to speak for themselves. A small community of people say exactly that, and the clearest
example in Britain is Charmaine D'Rozario-Saytch of Hove, on the south coast of England, who was
twenty-eight when the national press picked up her account in 2016.

Her claim is specific rather than vague. She describes herself as a human--extraterrestrial
hybrid of reptilian and lion-like lineage, says she was conceived through an off-world programme,
reports lifelong contact including having been shown other hybrid children she understands to be
her own, and says she can sense the presence of non-human intelligence. She appears in Barbara
Lamb and Miguel Mendon\c{c}a's \emph{Meet the Hybrids}, a 2015 collection of interviews with
twenty-six self-identified hybrids, and later on \emph{Ancient Aliens}. She has been consistent
for a decade, she is articulate, and --- this matters --- she is not obviously profiting.

So apply the chapter's own standard, evenly. What would settle it? Almost nothing about this
claim is beyond reach: a genome is cheap to sequence, and a person with non-terrestrial ancestry
would show it in every cell of their body. Reptilian and feline lineage in a human genome would
not be a subtle statistical signal --- it would be the most consequential laboratory result in
the history of biology, obtainable from saliva, for the price of a weekly shop. That test has
never been published for any self-identified hybrid, including this one. Where genetic testing
has been done on hybrid claimants and their alleged offspring, the results have come back human,
and where physical samples have been offered as anomalous --- the small mummified Atacama
specimen promoted as non-human in 2013 is the instructive case --- sequencing resolved them to
human beings with identifiable medical conditions.

I take no pleasure in that paragraph. The people in \emph{Meet the Hybrids} are not con artists
in any sense I can detect; several are plainly people for whom the identity has provided
meaning, community and an explanation for feeling permanently out of place, and that is a real
human need being met. But a claim about ancestry is a claim about DNA, and DNA is the one witness
in this entire book that cannot be led, cannot be regressed, and cannot confabulate. Twenty-six
interviews and one sequencing run were both available. Only the interviews were done.

% ---------------------------------------------------------------
\section{Missing 411}\label{sec:missing411}

\emph{Missing 411} is a series of books by David Paulides, a former police detective, cataloguing
disappearances in national parks and wilderness areas that he argues share anomalous
characteristics. He is careful --- deliberately and, I think, shrewdly --- not to state a cause.
The books present clusters, note recurring details, and leave the reader to supply an
explanation.

The recurring profile includes: disappearances near water or boulder fields; victims found with
clothing removed or neatly folded; children found improbable distances from where they vanished,
sometimes at elevation; dogs unable to pick up a scent; sudden weather at the time of
disappearance; German surnames; disappearances clustering geographically; and survivors with no
memory of the missing period.

Some of this is genuinely strange, and I do not want to be glib about a subject that involves
dead children. But the methodological problems are severe and they are of a specific kind worth
learning.

The first is the missing denominator, which is Section~\ref{sec:doubtdrivenskepticism}'s central complaint. Hundreds of
millions of visits are made to US national parks annually. Any set of features you select will
occur in some subset of cases. Without knowing the base rate --- how often do \emph{found}
missing persons have their clothes off, how often is there sudden weather in mountains, how
common are German surnames in the American population --- a list of coincidences is
unanalysable. Paulides does not provide base rates.

The second is that the anomalous details have well-documented mundane causes, and several are
signatures of specific known processes. Paradoxical undressing is a recognised feature of the
late stages of hypothermia: as the peripheral vasoconstriction fails, the victim feels a rush
of heat and removes clothing. It is documented in mountain rescue literature and in forensic
pathology, and it accounts for one of the eeriest items on the list. Terminal burrowing --- the
hypothermic victim's tendency to crawl into a confined space --- accounts for bodies found in
places that seem impossible to reach and that searchers walked past. Small children can and do
cover remarkable distances and climb, and they hide from searchers calling their names because
they have been taught to avoid strangers. Scent tracking fails routinely for entirely
understood reasons involving wind, water, temperature gradients and contamination. Weather in
mountains changes suddenly as a matter of course, which is why people die in mountains.

The third is a refusal-to-conclude problem. By declining to state a hypothesis, the books
become unfalsifiable while retaining all the rhetorical force of a claim. The National Park
Service has stated it does not maintain a comprehensive centralised missing-persons database of
the kind Paulides describes requesting, and his account of that refusal is presented as
suppression rather than as a records-management deficiency, which is what it is.

\begin{funfact}
A researcher named Kyle Polich attempted the obvious control: he compiled a comparison set of
wilderness disappearances that were resolved and checked how many exhibited the ``411''
profile markers. A great many did. When your anomaly profile is equally present in the solved
cases, you have not identified a pattern in disappearances. You have identified a pattern in
wilderness.
\end{funfact}

I include this section because Missing 411 is the most methodologically instructive body of work
currently circulating in the paranormal world. Everything in it is a real case. Every detail is
documented. And the whole is still worthless as evidence, for the single reason that nobody
counted what the ordinary rate would be. That is a more advanced failure than a hoax, and much
harder to spot.

% ---------------------------------------------------------------
\section{The Evidence for Abduction Cases}\label{sec:evidenceabductioncases}

I have spent five sections dismantling the abduction literature, so let me now do the thing
Section~\ref{sec:doubtdrivenskepticism} demands and lay out the strongest case for it, as fairly as I can.

\textbf{The consistency argument.} Thousands of accounts, from people of different ages,
countries, educational levels and cultures, share detailed features --- including details not
present in popular media at the time of the report. Some early cases predate the availability
of the script.

\textbf{The witness quality argument.} Abduction reporters are not, as a group, mentally ill.
This has been tested. A 1983 study by Ted Bloecher, Aphrodite Clamar and Budd Hopkins,
and more rigorously the work of psychologist Nicholas Spanos in 1993, found that experiencers
score within normal ranges on standard psychopathology measures. They are not psychotic, not
unusually fantasy-prone by clinical measures, and frequently reluctant --- many have lost
relationships and jobs over their accounts, and many kept quiet for years.

\textbf{The physical marks argument.} Some experiencers present scoop marks, straight-line
scars, and puncture wounds. A few cases involve claimed implants, and several have been
surgically removed and examined.

\textbf{The multiple-witness argument.} A handful of cases involve more than one person
reporting a shared abduction, and a very few involve independent observation of an object at
the same time and place --- the 1975 Travis Walton case of Section~\ref{sec:fireinthesky}, in which six logging crew members gave
consistent accounts of Walton being struck by a beam of light and reported his disappearance
to police the same night, is the strongest instance in the corpus.

\textbf{The authority argument.} John Mack, a tenured professor of psychiatry at Harvard Medical
School and a Pulitzer winner, investigated experiencers and concluded that conventional
psychiatric explanations were inadequate. Harvard convened a committee to review his work and,
after fourteen months, took no action against him.

Now the counter-argument, item by item, because the whole point of Chapter~\ref{ch:language} was that this is
the work.

Consistency: the script was culturally available from 1966 and dominant from 1987, and the
majority of the corpus was collected after both dates by a handful of investigators using
leading methods. The genuinely pre-script cases are few, and the small number that exist are
the only ones I regard as evidentially interesting.

\figwrap[r]{0.34}{Marks photographed on the bodies of several different self-described
abductees. The scars are genuine; what nobody has ever established is where they came from.}{}{fig:abdscars}
Witness quality: agreed, and it is important. But sanity is not accuracy. Perfectly healthy
people experience sleep paralysis, confabulate under hypnosis, and hold false memories with
complete conviction. Demonstrating that experiencers are not insane refutes a claim nobody
serious was making.

Physical marks: the marks are real. Their attribution is not. Human beings acquire
unexplained scars constantly and normally do not notice; once a person is looking for evidence,
they find their own body. Of the implants surgically removed and independently examined, none
has ever been shown to be of non-terrestrial manufacture --- the published analyses have found
glass, carbon fibre, and metals with terrestrial isotope ratios, which is decisive.
Section~\ref{sec:deconstructingpeerreview} explains why the isotope point is the one that settles it.

Multiple witnesses: Walton is the best case and it is still not clean. His account developed
over subsequent years, he had a documented prior interest in the subject, he failed one
polygraph in 1975 --- a fact suppressed for two decades --- and the crew's reports concern
what they saw from a truck, not what happened during the five days. As
Section~\ref{sec:fireinthesky} argues, those five days are the part with a physiological answer,
and the answer is not a spacecraft.

Mack: his integrity is not in question and his clinical observations about the effect of the
experience on his patients were careful and valuable. His conclusion, however, rested on
regression material, and no amount of professional standing repairs a compromised instrument.
Section~\ref{sec:deconstructingpeerreview} again: credentials are not a method.

\begin{funfact}
The cleanest experiment ever run on this subject is almost never mentioned. In 1993 Nicholas
Spanos and colleagues gave people who had never reported an abduction a hypnotic session with
leading suggestions --- and a substantial fraction produced detailed abduction narratives,
indistinguishable in structure and emotional intensity from those of self-identified
experiencers. That is the experiment. The narrative can be manufactured on demand, in a
laboratory, in an afternoon.
\end{funfact}

\subsubsection*{Roger Leir, and the question the implant debate never asks}

The implant claim deserves better treatment than the single clause I gave it above, because it is
the one strand of the abduction corpus that produces a physical object a laboratory can hold. The
man who produced most of those objects was Roger Leir, whose record is set out in
Section~\ref{sec:leir}: a podiatric surgeon who between 1995 and his death in 2014 carried out some
seventeen removals of foreign bodies from people who believed they had been abducted. My assessment
of his method there is unchanged --- no independent surgical oversight, no blinded pathology, no
chain of custody worth the name, and analyses commissioned rather than replicated --- and the
isotope verdict of Section~\ref{sec:chemistryufology} is unchanged with it. What Leir removed was
terrestrial. That much is settled and I am not going to soften it.

But settling the provenance does not settle the observation, and this is where almost everybody on
both sides of the argument stops one step too early. The debate has been conducted for thirty years
as a customs inspection: where was this object manufactured, and can we prove it came from
somewhere other than Earth. That framing guarantees the answer we already have, because a fragment
of terrestrial glass or carbon fibre in a human foot is an unremarkable object with an
unremarkable answer. Leir's own case files, read carefully, are not actually interesting because of
what the metal was. They are interesting because of what the tissue was doing.

\actualq{Do aliens implant human beings with tracking devices?}{Why does the body not reject these
objects as foreign material --- and why, in a number of the documented cases, does the surrounding
tissue fuse to them rather than wall them off?}

That is a physiological question, it is answerable, and it does not require a spacecraft to be
interesting. The normal response of mammalian tissue to an inert foreign body is well characterised
and it is hostile: acute inflammation, then a foreign-body giant-cell reaction, then a fibrous
capsule that isolates the intruder from the host, frequently with chronic pain, migration or
suppuration on top. This is why a splinter hurts, why surgical implants are engineered specifically
to negotiate with that reaction, and why an unsterile fragment lodged in soft tissue for years is
normally a clinical problem rather than a quiet passenger. \hi{Leir reported the opposite: objects
sitting in the tissue with no inflammatory infiltrate, no capsule of the expected kind, and a
nerve-rich, keratin-and-collagen sheath adherent to the object itself.} Several of his patients
reported no pain and had not known the object was there.

If that observation is real, it is the finding, and its explanation is almost certainly biological
rather than extraterrestrial. Innocuous fragments do sometimes become quiescent, particularly small,
smooth, biologically inert ones that provoke little surface reactivity; the body's tolerance of
embedded material varies enormously by site, by material and by patient; and a chronically
embedded object can become invested in scar tissue in ways that look like deliberate integration to
an eye already expecting design. Any of those would be a genuine contribution to knowledge. None of
them was ever pursued, because Leir's supporters wanted the object to be alien and his critics only
needed to prove it was not, and both sides therefore treated the histology as a footnote to the
metallurgy. \hi{The implant literature was answering a question about origin when the evidence it
held was about immunology.} That is the failure of language this whole book is about, and it cost
the field the one testable result the abduction corpus ever handed it.

My conclusion, and I want to state it in a way that does not insult anybody: the abduction
experience is real. People are genuinely experiencing something, it frightens them badly, and
it is not fraud and not madness. The evidence that anybody has been physically removed from
their bedroom by non-human entities is, after sixty years and thousands of cases, nil. Not
weak. Nil --- there is no artefact, no photograph, no medical finding, no missing-person
record, no independent corroboration of a single abduction event. A phenomenon this widespread,
if physical, would have left something behind by now.

% ---------------------------------------------------------------
\section{Introduction to Parapsychology}\label{sec:parapsych}

The chapter title promised parapsychology and here it is, and I owe you an explanation of why it is
bolted onto a chapter about abduction rather than given a chapter of its own.

The reason is that the two subjects share a nervous system. Abduction research and psychical
research both study a private experience that leaves no residue; both rely on testimony recovered
under conditions that can generate it; both have a century or more of accumulated case files and no
replicable result; and both, crucially, keep reporting the same secondary phenomena. Read enough
abduction literature and you will find telepathic communication, out-of-body episodes, precognitive
warnings, poltergeist activity in the household afterwards, and an insistence that the entities
operated on consciousness rather than on matter. John Mack noticed that. Jacques Vall\'ee built a
whole thesis on it. Whatever you think is happening, the phenomenology of abduction is closer to
the phenomenology of the s\'eance room than to the phenomenology of a flight test. So we should
look at what happened when serious people spent a hundred and forty years trying to nail the
s\'eance room down.

\subsection{Definitions, Properly Separated}

Parapsychology is the scientific study of claimed phenomena that, if real, cannot be accounted for
by currently understood physical mechanisms of information transfer or causation. The standard
taxonomy, which is worth learning because the field's precision about its own terms is genuinely
admirable:

\term{Psi} is the umbrella term. It divides into \term{ESP} --- extrasensory perception, meaning
information acquired without a sensory channel --- and \term{psychokinesis}, meaning influence on a
physical system without a physical channel. ESP subdivides into \term{telepathy} (mind to mind),
\term{clairvoyance} (mind to remote object or event), and \term{precognition} (information from the
future). Psychokinesis subdivides into \term{macro-PK}, the visible kind that Uri Geller performed
and that has never once survived controlled conditions, and \term{micro-PK}, the statistical kind:
tiny deviations in random number generators, which is where the modern field actually lives.

Separately there is \term{survival research} --- mediumship, apparitions, reincarnation claims,
near-death experiences --- which asks whether consciousness persists after death. Section~\ref{sec:whatghost} and
Section~\ref{sec:afterlife} handle that.

\subsection{A Very Short History}

The Society for Psychical Research was founded in London in 1882, and its founding membership is
the part that surprises people: Henry Sidgwick, the Cambridge moral philosopher, was its first
president. Later presidents and members included the physicists William Barrett, Oliver Lodge and
Lord Rayleigh, the psychologist William James, the philosopher Henri Bergson, and the Prime
Minister Arthur Balfour. This was not a fringe operation. It was the Victorian intellectual
establishment attempting to apply experimental method to spiritualism, and its early achievement was
largely negative and largely excellent: it exposed fraudulent mediums in detail, including some its
own members badly wanted to be genuine.

J.\,B.\ Rhine moved the subject into the laboratory at Duke University in the 1930s, coined the term
``extrasensory perception'', and introduced the Zener card deck and statistical evaluation. Rhine's
contribution was methodological --- he turned anecdote into a countable experiment --- and his
results did not replicate outside his own laboratory, and in at least one case involved a
collaborator, Walter Levy, caught falsifying data in 1974.

The cast of the government period deserves naming, because their biographies are themselves an
argument. Hal Puthoff, a laser physicist with a Stanford doctorate in electrical engineering,
built the protocols and ran the paranormal research programme at the Stanford Research Institute
through the 1970s. Ingo Swann, an artist and self-described psychic who claimed abilities from
early childhood, is generally credited with co-developing the coordinate remote-viewing method.
Joseph McMoneagle, who enlisted in 1964 and came to the work through Army intelligence, was
designated Remote Viewer No.~1 when the operational unit stood up in 1977 and is the only one of
the three who was ever formally decorated for the output. There were roughly three dozen
military volunteers through the programme's life. And here is the detail that nobody in the
field's promotional literature volunteers: Puthoff and Swann were both, at the relevant period,
members of the Church of Scientology, an organisation whose doctrine already asserted that the
abilities being tested were real and recoverable by training. Several of the principals met
either at SRI or through the church.

I raise it not as a smear but as a methodological fact of exactly the kind Section~\ref{sec:confirmationbias} told you
to look for. When the people designing the experiment, the people being tested, and the people
scoring the transcripts share a prior religious commitment to the phenomenon existing, the
resulting protocol will leak. That is not a claim about anyone's honesty; it is a claim about how
blinding fails. It also explains the programme's peculiar signature --- results that were
statistically odd in-house and vanished when outsiders ran the sessions. And it is why the
lineage matters: the remote-viewing work sat downstream of MKULTRA, which had already spent a
decade demonstrating that the American security establishment would fund almost any hypothesis
about the mind, and had already produced an institutional habit of running human research
without controls that would survive external review.

Then came the government period, which is the part relevant to the rest of this book. Between 1972
and 1995 the CIA and the Defense Intelligence Agency funded remote-viewing research at the Stanford
Research Institute and later at SAIC, under a sequence of code names ending with STARGATE. Real
money, roughly \$20 million, real classification, real operational tasking. The programme was
terminated in 1995 after an external review commissioned by the CIA --- Jessica Utts, a
statistician sympathetic to the field, and Ray Hyman, a psychologist who was not --- produced a
remarkable pair of reports. Utts concluded that a statistically significant effect had been
demonstrated. Hyman agreed that the statistics were anomalous but held that methodological
flaws and the absence of any mechanism made a paranormal conclusion unwarranted. Both agreed the
programme had never produced intelligence of operational use. Read those two reports side by side;
they are the single best education available in what a genuinely disputed borderline result looks
like, as opposed to the fake disputes that fill most of this field.

\begin{funfact}
The STARGATE files were declassified and are on the CIA's own public reading room, which means the
sentence ``the CIA studied psychic phenomena'' is entirely true and is also the setup for the most
misleading move in the paranormal repertoire. Yes, they studied it. For twenty-three years. Then
they read the assessment and shut it down. The funding is evidence that the question was taken
seriously; the termination is evidence about the answer. People quote the first half.
\end{funfact}

\subsection{Where the Evidence Actually Stands}

I want to be careful here, because parapsychology is the one place in this book where the honest
sceptical position is more interesting than a flat dismissal.

\figwrap[l]{0.38}{Ganzfeld apparatus: halved ping-pong balls, red light and white noise, producing uniform sensory input.}{https://commons.wikimedia.org/wiki/File:Ganzfeld_experiment_setup.png}{fig:ganzfeld}
The \textbf{ganzfeld} telepathy experiment is the field's best paradigm. A receiver in mild sensory isolation (see Figure~\ref{fig:ganzfeld}) attempts to identify which of four images a distant sender viewed; chance is 25\%.
Meta-analyses across hundreds of sessions have repeatedly reported hit rates in the low-to-mid 30s,
which is small but not trivial. Charles Honorton and Ray Hyman jointly published a ``joint
communique'' in 1986 agreeing on stricter protocol standards --- a genuinely honourable episode in
the history of scientific disagreement. Under the tightened protocols the effect persisted in some
labs and vanished in others.

The \textbf{random number generator} experiments, particularly the Princeton Engineering Anomalies
Research programme run out of Princeton's engineering school from 1979 to 2007, reported deviations
from expectation on the order of one part in ten thousand, accumulated across tens of millions of
trials. PEAR's own conclusion after twenty-eight years was that the effect was real but too small
and too unstable to characterise.

And in 2011 Daryl Bem published nine experiments in the \emph{Journal of Personality and Social
Psychology} reporting precognition --- that participants' responses were influenced by stimuli
presented afterwards. This is the most consequential paper in the history of parapsychology, and not
for the reason Bem intended. It used entirely standard social-psychology methods, passed peer review
at a mainstream journal, failed to replicate, and in doing so helped trigger the replication crisis
in psychology as a whole. Bem's paper is now taught in methods courses as the reductio: if these
normal, accepted, everyday research practices can produce a significant result for something
impossible, then those practices cannot be trusted to establish anything.

That is the state of play. A century and a half of work, a serious methodological tradition, and a
residual effect that is always small, always near the detection threshold, never gets larger with
better instruments, and shrinks as controls tighten. Compare that with any real phenomenon: the
effect size for aspirin on headache, or the charge on an electron, does not diminish when you
improve the apparatus. Real effects get \emph{clearer}. This is the diagnostic test, and psi has
been failing it in the same direction for a hundred years.

\pullquote[Ray Hyman]{``The case for psychic functioning seems better than it ever has been \dots\ I
also have to admit that I do not have a ready explanation for these observed effects.''}

Why include all this in a UFO textbook? Because it is the closest thing we have to a completed
experiment on our own discipline. Parapsychology did everything Ufology is constantly told to do.
It got into universities. It adopted statistics, controls, blinding, pre-registration, and
meta-analysis. It obtained government funding and classified operational trials. It attracted
physicists and Nobel laureates. It ran for a hundred and forty years. And it still cannot
demonstrate its central claim on demand. That should be sobering for anyone who believes the only
thing standing between Ufology and a definitive answer is respectability and a research budget. It
is possible to do this properly, at length, with money, and still find nothing --- and it is
possible that finding nothing is the correct outcome, honestly arrived at.

\begin{srcnotes}
\item Utts, J.\ (1995) and Hyman, R.\ (1995). Evaluations of the government remote-viewing
programme, prepared for the American Institutes for Research review of STARGATE.
\item Bem, D.\ J.\ (2011). Feeling the Future. \emph{Journal of Personality and Social
Psychology}, 100(3).
\item Hyman, R.\ \& Honorton, C.\ (1986). A Joint Communiqu\'e: The Psi Ganzfeld Controversy.
\emph{Journal of Parapsychology}, 50.
\item Jahn, R.\ G.\ \& Dunne, B.\ J., Princeton Engineering Anomalies Research, final reports,
2007.
\end{srcnotes}

% ---------------------------------------------------------------
\section{Ghost Hunting}\label{sec:ghosthunting}

A reader is entitled to stop here and ask what ghost hunting is doing in a textbook about
unidentified flying objects. The answer is one instrument, and it is worth stating before the
section begins rather than after it. Ghost hunting is the field that decided, and then taught
everybody else, that \hi{electromagnetism is the signature of the paranormal}. Every kit sold to
every investigator in this trade is built around a meter that reads a field strength in a room, on
the premise that a reading means a presence. That premise did not stay in the haunted house. It
was borrowed wholesale by the rest of the anomalous literature --- by abduction researchers, by
investigators of landing traces, and most consequentially by the best-funded field programme in
modern ufology, the one examined in Section~\ref{sec:skinwalker}, where a bulb lit by a car
battery in conductive Utah soil was read on air as something beneath the ground announcing itself.
None of that later material can be assessed by a reader who has not first watched the equation
between field and presence being made here, in the trade that made it. So the order is deliberate.
Learn the meter with ghosts, where the stakes are low and the footage is abundant, and by the time
the same meter is pointed at a ranch in the Uinta Basin you will already know exactly what it does
and does not measure.

\subsection{Who Sam and Colby Are}\label{sec:whosamcolby}
Before the definitions, the people. Much of the field material in this chapter is drawn from
two American video makers, and a reader who has never heard of them deserves to know who
they are before their footage is used as evidence of anything.

\figwrap[r]{0.42}{The \emph{Queen Mary} at Long Beach, California --- one of the locations Sam
and Colby have filmed in repeatedly, and a useful reminder that these are real, named,
publicly accessible buildings rather than anonymous sets. No free-licence photograph of
Golbach or Brock themselves exists at the time of writing: every available portrait is
commercially licensed or of unverifiable provenance, so the site is shown instead of the
investigators.}{https://commons.wikimedia.org/wiki/File:RMS_Queen_Mary_Long_Beach_January_2011_view.jpg}{fig:samcolby}
Sam Golbach and Colby Brock are both from Kansas, and they met as teenagers at the same high
school in the Kansas City suburbs. They began posting together in the mid-2010s on Vine, the
short-video platform of the day, and when that platform closed they moved to long-form video
on YouTube. The channel they built there, \emph{Sam and Colby}, has since grown past fourteen
million subscribers, and the surrounding brand --- \emph{XPLR} --- covers merchandise, a
subscription app and a touring live presence. By any ordinary measure of reach they are among
the most widely watched paranormal investigators alive, and they reach an audience roughly the
size of a mid-sized European country without ever appearing on television.

What they actually do is stay overnight, on camera, in named locations with documented
histories: the \emph{Queen Mary} at Long Beach, the Sallie House in Atchison, the Villisca axe
murder house in Iowa, and dozens of hotels, prisons, asylums and private homes besides. The
format is consistent --- multiple nights in the same building, multiple cameras running,
long stretches of nothing, and a small number of episodes in which something happens that
they cannot explain and do not pretend to. They also film the Estes Method sessions discussed
later in this chapter, and they film them in enough volume that the sessions can be compared
against each other rather than treated as one-off anecdotes.

None of that makes them scientists, and they have never claimed to be. It does make them
prolific, consistent and unusually transparent recorders of a phenomenon that academic
research has largely refused to go anywhere near, and that is the reason they appear in a book
about evidence. The relevant question is not whether two men in their twenties with a camera
rig are credentialled. It is whether the material they have accumulated is the sort of thing a
serious investigator can work with. The rest of this chapter argues that in several specific
respects it is, and is honest about the respects in which it is not.

\subsection{What is a GHOST?}\label{sec:whatghost}
You may reasonably ask what ghosts are doing in a book about UFOs. Two answers. First,
because the field itself keeps mixing them --- a substantial minority of UFO cases contain
poltergeist activity, apparitions and precognition, and Skinwalker Ranch in
Section~\ref{sec:skinwalker} is unintelligible without acknowledging it. Second, and more
importantly, because ghost hunting is the perfect control group. It is a field with more
enthusiasts, more equipment, and better site access than Ufology has ever had, and it has
produced nothing in a century and a half. Studying how that happened is cheaper than repeating
it.

There is a third answer, and it is the practical one, so it is worth stating before the
equipment section rather than after it. \hi{Ghost hunting is where the reader learns that in the
paranormal trades electromagnetism \emph{is} the evidence.} That premise is not a detail of one
hobby. It is the shared instrument of the entire field: the EMF meter is the first thing bought,
the first thing switched on, and the thing whose needle is read out loud as though a field
strength were a presence. A reader who has watched that inference being made here, on a subject
where nobody feels obliged to be careful, is equipped to recognise it later on ground where the
stakes are higher --- at Skinwalker Ranch in Section~\ref{sec:skinwalker}, where a television
crew lights a bulb through fifty feet of saline Utah soil and reads the result as something
beneath the field. The same meter, the same reasoning, a different genre. You cannot audit that
argument in a chapter on UFOs unless you have first seen where it was learned.

``Ghost'' is at least five distinct claims wearing one word, and separating them is most of the
work:

\textbf{The survival hypothesis.} A ghost is the persisting consciousness of a dead person.
This requires consciousness to be separable from the brain, which is discussed properly in
Section~\ref{sec:afterlife}.

\textbf{The recording hypothesis.} The ``stone tape'' idea: environments somehow store and
replay events. No proposed physical mechanism exists, and stone has no known information
storage capacity of the required kind.

\textbf{The apparitional experience.} A perceptual event occurring in the observer.
Well-documented, common, and requiring no external entity at all.

\textbf{Poltergeist activity.} Physical disturbance --- objects moving, sounds, electrical
failures --- usually centred on a person rather than a place, and disproportionately on an
adolescent, which the older literature noticed and drew exactly the wrong conclusion from.

\textbf{Interdimensional or non-human entities.} Vall\'ee's ultraterrestrial idea, and
Section~\ref{sec:shadow}'s subject.

The conventional explanations for hauntings are unusually well established, and two are worth
knowing in detail.

Carbon monoxide poisoning produces, at sub-lethal chronic doses: headache, nausea, confusion,
dread, auditory hallucination and visual apparition. There is a celebrated 1921 case ---
reported in the \emph{American Journal of Ophthalmology} --- of a family who moved into a house
and experienced footsteps, voices, and figures at the bedside, all of which resolved when a
faulty furnace was found. Any haunting confined to one building with a heating system should be
investigated with a gas detector before an EMF meter.

Infrasound --- sound below about 20\,Hz, inaudible but perceptible --- produces anxiety,
chills, sensations of presence, and, near 18--19\,Hz, visual disturbance, because that
frequency is close to the resonant frequency of the human eyeball. Vic Tandy, an engineer,
identified this in 1998 after tracing a ``haunted'' laboratory to a newly installed extractor
fan; when the fan was fixed, the ghost left. Old buildings with large fans, plumbing, or
wind-driven structural resonance generate infrasound routinely.

Add to those: mould, particularly \emph{Stachybotrys}, associated with cognitive symptoms in
some studies; sleep paralysis, as in Section~\ref{sec:evidenceabductioncases}, which accounts for the enormous class of
bedroom apparitions; and pareidolia and expectation, which account for the rest. A person told
a house is haunted will experience a haunted house.

\begin{funfact}
The suggestion effect has been measured directly. Richard Wiseman ran studies at Hampton Court
Palace and the Edinburgh Vaults in which visitors were led through locations and asked to
report unusual experiences. Reports clustered strongly in areas with prior haunting
reputations --- and also correlated with measurable environmental variables: air movement,
temperature variation, lighting, and ceiling height. People reliably feel watched in places
with certain physical properties. The ghost is in the architecture.
\end{funfact}

\subsection{The Instruments, One at a Time}\label{sec:ghostinstruments}
Modern ghost hunting is a methodological catastrophe conducted with great sincerity, and I want
to go through the instruments one at a time, because every criticism here transfers directly to
UFO fieldwork. The UFO investigator's kit in Section~\ref{sec:equipmenttools} --- Geiger counter,
spectrometer, thermal camera, magnetometer --- was assembled to answer questions about physical
traces. The ghost hunter's kit was assembled to answer questions about intent. That is the whole
difference, and it is fatal.

I am going to use one source for the field practice, and I want to say why. Sam Golbach and Colby
Brock --- Sam and Colby --- are two American YouTubers with an audience north of fourteen million,
which makes them, by a wide margin, the most-watched paranormal investigators who have ever
lived. More people have watched them run a spirit box session in a night than have read every
book cited in this chapter combined. Their kit is the standard kit, bought from the same handful
of suppliers as everyone else's, and their protocol is the standard protocol. So when I describe
what the field actually does at two in the morning, I am describing what they do, and I am not
doing it to sneer. They are conspicuously not frauds: they film the boring stretches, they say
so when nothing happens, they leave in the moments where one of them talks himself into a
fright, and they are visibly braver about entering condemned buildings than I have ever been
about anything. The problem with their investigations is not honesty. It is that the instruments
they have been sold cannot do the job the packaging implies, and nobody in the supply chain has
any commercial reason to tell them.

\figwrap[r]{0.32}{A generic handheld EMF meter of the type sold to investigators. It is an honest instrument measuring a real field --- usually one produced by the building's wiring or somebody's phone.}{}{fig:emfmeter}
\textbf{The EMF meter.} Used to detect ``spirit energy''. It detects electromagnetic fields,
of which a house is full: wiring, appliances, phones, the meter operator's own equipment, and
the investigator's radio. There has never been any demonstrated link between EMF readings and
apparitional experience --- and note the irony, since if there is a real effect here it runs
the other way: strong fields may cause the experience, which would make the meter a detector
of causes rather than of ghosts. Watch any investigation carefully and you will see the meter
spike as somebody's phone hunts for signal in a basement, which is a real detection of a real
field produced by a real transmitter carried by a real person standing four feet away.

\figwrap[l]{0.32}{A single-axis EMF gauge and thermocouple probe of the general type combined in the meter described here. Competent hardware; the difficulty is that no reading has ever been assigned a meaning.}{}{fig:melmeter}
\textbf{The Mel Meter (8704R and relatives).} Worth singling out, because its history is the most
human thing in the entire industry. Gary Galka, an instrumentation engineer, lost his
daughter Melissa in a car accident in 2004; the meter is named for her --- \hi{MEL, her initials}
--- and he built it because he believed she was trying to reach him. What he actually built was a
competent piece of hardware: a combined single-axis EMF gauge and thermocouple thermometer with a
red-lit display, so the operator's night vision survives the reading, and both values on one
screen so the two can be logged together. Every design decision solves a genuine field problem.
And none of them addresses the only problem that matters, which is that nobody has established
what reading would indicate a ghost. Better ergonomics on an uninterpretable measurement produce
uninterpretable measurements more comfortably.

\figwrap[r]{0.32}{A theremin --- the capacitive-antenna instrument whose operating principle the REM-POD borrows. A radiating aerial senses changes in its own capacitance as a body approaches.}{}{fig:rempod}
\textbf{The REM-POD (EMT 2.0, with ATDD).} This one deserves a moment because it is the cleverest
object in the kit and the most badly reasoned about. A REM-POD does not passively listen. It
\emph{radiates} --- a telescopic antenna generates its own weak electromagnetic field, and the
device monitors that field for disturbance, sounding off and lighting up when something alters the
antenna's capacitance. The EMT~2.0 adds ATDD, ambient temperature deviation detection, which
alarms on a swing of a degree or two either way. As a proximity sensor it works exactly as
advertised. The trouble is what counts as proximity: a hand, a sleeve, a damp coat, a moth, a
shift in humidity, a mobile phone, a two-way radio, or a second REM-POD placed helpfully nearby.
Investigators know this and set the pod down as a tripwire, then stand around it talking. And the
temperature alarm, in an unheated building at night, is triggered by the thing unheated buildings
do at night, which is cool unevenly as air moves through them. \hi{An alarm that fires when a
person approaches is not evidence when a person approaches.}

\figwrap[l]{0.30}{A non-contact infrared thermometer. It reads the surface inside a cone that widens with distance, and never reads air temperature at all.}{}{fig:irtherm}
\textbf{The infrared thermometer.} A non-contact IR thermometer reads the surface temperature of
whatever is inside its cone, and the cone widens with distance --- at ten feet, a device with a
12:1 optical ratio is averaging a patch of wall the size of a dinner plate. It does not measure air
temperature at all, which is the single most consequential misunderstanding in the field, because
``cold spot'' claims are claims about air. Point it at a mirror and you get the temperature of
whatever the mirror is reflecting. Point it at a window and you get the outside. Sweep it across
a room and you will find a five-degree variation in any room ever built. The instrument is honest;
the inference is not.

\figwrap[r]{0.30}{A passive infrared motion sensor. Instruments in this category could settle an argument --- if anyone ran them in an empty, sealed room.}{}{fig:motionsens}
\textbf{Motion, contact and trigger devices.} Laser grid pens, PIR motion alarms, pressure mats,
door contacts, and the trigger object --- a ball, a coin, a toy left on a chalk outline. This is
the most promising category and the most consistently wasted. A pressure mat and a door contact
are unambiguous: something either weighs on the mat or it does not. Used properly --- sealed
room, everyone outside, continuous timecoded camera on the object, chalk outline photographed
before and after --- they are capable of producing a result that would survive an argument. Used
as they normally are, with four people in the room and the object handled between takes, they
produce nothing at all. The field owns the equipment to run a controlled experiment and reliably
chooses not to run one, because a sealed room with nobody in it makes very poor television.

\figwrap[l]{0.30}{A twist-cap torch of the familiar kind. Loosen the cap to the edge of contact and thermal expansion will do the rest, on roughly the timescale at which questions are asked.}{}{fig:torch}
\textbf{Flashlights.} The twist-cap torch, unscrewed to the point where the contact is marginal,
asked to turn itself on for yes. This is not a paranormal detector, it is a demonstration of
thermal expansion: a Maglite-style barrel and bulb assembly warms, cools, and expands and
contracts across exactly the tolerance the operator has set it to, and the contact makes and
breaks on a timescale of tens of seconds, which happens to be the timescale on which people ask
questions. Set two of them up and leave the room, and they will carry on holding a conversation
with nobody.

That much is mechanism, and I want to be careful not to overstate what it settles. Hours of
flashlight work exist on camera --- Sam and Colby alone have filmed a great deal of it, and the
torches in that footage behave exactly as the people holding the session describe: they come on
when asked and they stay off when asked to stay off. So the interesting question is not the one
the sceptical literature usually reaches for. Whether an intelligence can move a physical object
is not seriously in dispute as an open question; a hundred and fifty years of poltergeist reports,
some of them investigated by people who badly wanted a mundane answer, have kept it on the table.
A torch barrel is a trivial load compared with the furniture in the Rosenheim or Enfield files.
The question a twist-cap torch actually raises is much narrower and much sharper: \emph{why does
the contact make and break in time with intelligent yes/no questioning?} Thermal cycling explains
switching. It does not explain answering.

\actualq[question]{whether something discarnate could switch a torch on --- if a poltergeist can
throw a chair across a kitchen, a marginal contact is nothing.}{why a marginal contact should
keep time with the semantic content of a question, rather than with the thermal cycle of the
barrel.}

And that is a question with a clean answer available, which is what makes the field's refusal to
reach for it so telling. Set two torches to the same tolerance. Put one in a sealed empty room
with a locked-off timecoded camera and nobody within earshot, run the other in the session, and
log every make and break on both. Then have the question timings and the switching timings scored
separately, by people who cannot see the other stream, and check whether hits cluster on questions
at a rate the unattended control does not reach. That is a weekend of work with equipment the
community already owns. Until it is run, every torch session is measuring the operator's tolerance
setting and the room's temperature, and the yes/no structure is being supplied entirely by the
person asking. \hi{The mechanism is not the mystery; the timing is, and nobody has measured the
timing.}

\figwrap[r]{0.32}{A full-spectrum converted camera of the type used for infrared work. On-axis flash plus near-field dust is the whole of the orb phenomenon.}{}{fig:fullspec}
\textbf{The full-spectrum or infrared camera.} Infrared photography captures reflected IR.
Dust, insects, moisture droplets and lens flare are all rendered as bright ``orbs''. The orb
phenomenon appeared in the literature at precisely the moment consumer digital cameras with
built-in flash became common, for the entirely understood optical reason that a flash close to
the lens axis illuminates near-field dust. There is no orb literature from the era of
large-format film with off-axis flash.

\figwrap[l]{0.32}{A thermal imaging camera. Genuinely useful for finding draughts, insulation gaps and pipework --- which is to say, for finding the mundane explanation.}{}{fig:thermphone}
\textbf{Thermal imaging.} A thermal camera shows surface temperature. It is genuinely useful
for finding drafts, insulation gaps, and plumbing --- which is to say, for finding the
conventional cause. It is routinely used instead to photograph ``cold spots'' whose measured
magnitude is within instrument tolerance.

\figwrap[r]{0.30}{A handheld digital audio recorder of the type used for EVP work. The recorder is not the problem; the listener with a transcript in hand is.}{}{fig:evprec}
\textbf{Audio recording and EVP.} Electronic voice phenomena --- voices claimed to appear in
recordings. The mechanism is auditory pareidolia, and it is one of the strongest perceptual
effects known. The human auditory system is aggressively tuned to extract speech from noise;
given random audio and a suggestion of what to listen for, most people will hear it clearly,
and once heard cannot un-hear it. Broadband noise, radio bleed, and digital compression
artefacts supply the raw material. The decisive test --- play the clip to listeners with no
prior suggestion and no transcript, and see whether they independently agree on the words ---
is essentially never performed.

\figwrap[l]{0.30}{A portable AM/FM receiver of the type a spirit box is built from. Sweep the band fast enough and a suggestible listener supplies the words.}{}{fig:spiritbox}
\textbf{The spirit box.} A device that sweeps AM/FM frequencies rapidly, producing fragments
of broadcast audio. This is a random word generator connected to a suggestible listener. It
cannot fail to produce apparent responses, which is exactly why it must not be used --- with one
qualification, which is interesting enough to have earned its own subsection below.

\ufowrapclear
\begin{funfact}
Every ghost-hunting technique shares one design flaw: not a single one has a null condition. No
instrument in the standard kit can produce a result meaning ``nothing is here''. An EMF meter
reading zero means the ghost is inactive; a reading means the ghost is present. Silence on the
recorder means try again. This is the Section~\ref{sec:claimsmadeslipperyslope} problem in hardware. A detector that cannot
register an absence is not a detector, it is a prop --- and if you have ever wondered why a
hundred and fifty years of enthusiastic effort has produced no result, that is the reason, and
it is entirely sufficient.
\end{funfact}

\ufowrapclear
\subsection{The Estes Method}\label{sec:estes}
The Estes Method is the only thing in modern ghost hunting that I think is genuinely worth
arguing about, and the reason is that whoever designed it accidentally solved the problem I have
just spent five pages complaining about.

The origin is prosaic. Around 2016, at the Stanley Hotel in Estes Park, Colorado --- the hotel
that gave Stephen King \emph{The Shining} --- three investigators named Karl Pfeiffer, Connor
Randall and Michelle Tate were running a spirit box and hit the obvious objection: the listener
can hear the questions, so the listener knows what answer would be impressive. Their fix was to
remove the listener from the room's information. One person, the receiver, is blindfolded, wearing
over-ear headphones fed directly from the sweeping spirit box, loud enough to mask all external
sound. A second person, out of earshot in the ordinary sense, asks questions aloud and writes
down both the questions and whatever the receiver says. The receiver speaks whatever words seem
to surface. Nobody tells them anything until the session ends and the transcript is compared.

Look at what that is. \hi{That is a blinded protocol.} Not a good one --- the questioner is in
the same room, headphone isolation is imperfect, the transcript is produced by an interested
party, and nobody randomises or counterbalances anything --- but it is a real attempt at
information control, invented by hobbyists, for the correct reason, and it is more methodological
self-awareness than the rest of the kit combined. If you want to know why the Estes Method spread
through the field within a few years while the Ovilus quietly died, that is why: it produces
results that feel like they came from somewhere the receiver could not have got them, because
some of the routes have actually been closed off.

Sam and Colby's sessions are the reason most people have heard of it, and they are worth watching
with a notebook. The hits, when they come, are of a very particular kind. Names. Numbers a
questioner was thinking of. Answers that respond to the \emph{sense} of a question rather than
its words. Occasionally something about a person in the room that the receiver would have no
obvious way to know. It is compelling on screen and it is compelling for a reason, and I want to
state plainly what I think that reason is, because it is not ghosts.

\hi{The Estes Method is remote viewing with extra steps, and the field that invented it does not
know that.} Compare it, line by line, with the Stargate protocol described earlier in this
chapter. A subject is placed in sensory-restricted conditions --- Ganzfeld used half ping-pong
balls over the eyes and white noise in the ears; Estes uses a blindfold and a sweeping radio. The
subject is asked to report whatever content arrives unbidden, without filtering it for sense.
A second party holds the target information and is kept, at least nominally, from cueing. The
output is a stream of fragmentary impressions --- words, shapes, names --- that is scored
afterwards against the target by comparison. Change the vocabulary from ``spirit'' to
``target'' and the two procedures are the same experiment. Estes Park reinvented the Ganzfeld,
forty years late, in a hotel corridor, and put a radio in it.

Which means the interesting question is not whether the spirits are talking. It is whether the
Estes Method performs better than chance --- and that question already has a large,
well-funded, thoroughly argued literature attached to it, which the ghost-hunting community has
never read and which concluded, after a century of trying, that effects of this kind hover
indistinguishably close to zero. If the Estes Method really is Ganzfeld-with-static, then the
honest prediction is that it will produce exactly the same result: striking individual sessions,
no reliable aggregate signal, and a debate that outlives everyone in it.

But --- and this is why I have given it a subsection rather than a paragraph --- it is testable,
cheaply, tomorrow, with equipment the field already owns. Put the questioner in a different
room on an intercom. Randomise the questions from a sealed list the questioner has not seen.
Run sessions with no questioner at all as a control. Have the transcript scored by someone
blind to which session was which. Nobody has done this properly, and the people best placed to
do it have fourteen million subscribers and a production budget. I would watch that video, and
I suspect it would be the most useful hour the paranormal community has produced in fifty years.
The full comparison --- Ganzfeld scoring methods applied to Estes transcripts, with the numbers
--- is more than this edition has room for, and it is top of my list for the second.

\begin{srcnotes}
\item Randall, C.\ J.\ (2019--2023). Interviews on the origin of the Estes Method at the Stanley
Hotel, including \emph{Haunted Road} with Amy Bruni and \emph{Bloody Disgusting}, March 2023.
\item Pfeiffer, K., Randall, C.\ J.\ \& Tate, M.\ \emph{Spirits of the Stanley} (2019), in which
the blindfold-and-headphones configuration is documented on film.
\item Sam and Colby (Golbach, S.\ \& Brock, C.). Paranormal investigation channel, YouTube; Estes
Method sessions across the Conjuring House, Sallie House and Winchester Mystery House series.
\item Sam and Colby (Golbach, S.\ \& Brock, C.). Flashlight sessions, various episodes --- the
largest filmed body of twist-cap torch work, and the reason the claim that the technique is
never recorded cannot be made. No unattended-control torch appears in any of it.
\item Bender, H.\ (1968--1969). The Rosenheim poltergeist investigation, and Playfair, G.\ L.\ \&
Grosse, M.\ (1977--1979), the Enfield case files --- object displacement reported under
investigation, which is why the switching of a torch is not itself the difficulty.
\item Honorton, C.\ (1985) and Bem, D.\ J.\ \& Honorton, C.\ (1994). The Ganzfeld protocol and
its meta-analyses --- the sensory-isolation design that the Estes Method independently
reinvented.
\end{srcnotes}

One last observation, and it applies to this whole chapter. There is a version of this research
that is done properly: it is called parapsychology, it has occupied laboratories at Duke,
Princeton, Edinburgh and elsewhere for the better part of a century, it uses randomised
protocols, blinding, pre-registration and meta-analysis --- and its results, after all that
rigour, hover indistinguishably close to chance. That is a real answer, arrived at honestly, and
it is far more interesting than any orb photograph. If a field with laboratory conditions,
funding, and unlimited access to its subject matter found nothing, then a field working with
handheld meters in the dark at two in the morning is not going to find it either.
